\subsection{Sáng tạo luật chơi mới - Biến thể Saboteur với  ngày đêm }
Biến thể mới giữ nguyên tinh thần bluffing và phản bội của Saboteur gốc nhưng bổ sung cơ chế pha \textbf{ngày} - \textbf{đêm}, thời gian giới hạn, tranh luận, bỏ phiếu và một số thẻ hành động mới mạnh mẽ hơn. Trò chơi vẫn chia hai phe: \textbf{Thiện (Gold-Diggers - Chó)} và \textbf{Ác (Saboteurs - Mèo)}, sử dụng bộ bài gồm thẻ đường (path cards) và thẻ hành động (action cards). Mục tiêu thắng thua giữ nguyên như phiên bản gốc.
\begin{enumerate}
\item \textbf{Chuẩn bị trước khi bắt đầu game}
\begin{itemize}
\item Chia vai trò (role cards) theo bảng phân bổ của phiên bản gốc (xem bảng \ref{tab:total-roles} - số lượng Saboteurs và Gold-Diggers theo tổng số người chơi).
\item \textbf{Thêm số lượng thẻ đường}: Sử dụng toàn bộ thẻ đường của bộ gốc, nhưng \textbf{ngẫu nhiên loại bỏ 10 thẻ đường} trước khi xào bài và để không ai biết chính xác số lượng cũng như loại thẻ đường nào còn lại trong bộ bài.
\item Chia bài cho mỗi người chơi như cũ theo bảng \ref{tab:cards-per-player} 
\item Xào phần bài còn lại thành chồng bài úp (draw pile).
\item Bố trí bàn chơi giống phiên bản gốc (xem hình \ref{fig:cards-and-board}).
\end{itemize}
\item \textbf{Cấu trúc một round}
\\
Mỗi round được chia thành hai pha chính: \textbf{Đêm} và \textbf{Ngày}.
\\
\textbf{Đêm (Night Phase):}
\begin{itemize}
\item Màn hình/bản đồ tối lại, người chơi \textbf{không nhìn thấy bất kỳ thay đổi nào trên bản đồ} ngoại trừ thông tin lượt hiện tại.
\item Thứ tự lượt chơi trong đêm được \textbf{xếp ngẫu nhiên} (không theo chiều kim đồng hồ như bản gốc).
\item Mỗi người chơi lần lượt thực hiện lượt:
\begin{itemize}
\item Bản đồ tạm thời hiện ra cho riêng người đó.
\item Trong \textbf{15--30 giây} (có thể tùy chỉnh), người chơi thực hiện một trong các hành động quen thuộc:
\begin{enumerate}
\item Đặt thẻ đường (path card) – phải tuân thủ đầy đủ quy tắc kết nối và xoay thẻ như bản gốc.
\item Chơi thẻ hành động (action card).
\item Bỏ bài (discard).
\end{enumerate}
\item Sau khi hết thời gian hoặc xác nhận hành động, bản đồ lại tối đi với người chơi đó.
\item Bốc 1 lá bài từ chồng bài úp (nếu còn).
\end{itemize}
\item Khi tất cả người chơi đã hoàn thành lượt trong đêm, pha Đêm kết thúc và chuyển sang pha Ngày.
\end{itemize}



\textbf{Ngày (Day Phase):}
\begin{itemize}
\item Bản đồ được chiếu sáng hoàn toàn, mọi người chơi cùng nhìn thấy toàn bộ thay đổi đã xảy ra trong pha Đêm.
\item Người chơi có \textbf{60 giây} (có thể tùy chỉnh) để \textbf{trò chuyện, tranh luận, biện minh} và đọc vị nhau.
\item Cuối pha Ngày, thực hiện \textbf{bỏ phiếu}:
\begin{itemize}
\item Có thể vote \textbf{skip} (không loại ai) hoặc vote một người chơi cụ thể.
\item Người bị vote nhiều phiếu nhất sẽ bị \textbf{cấm chơi 1 round tiếp theo} (không được hành động trong pha Đêm của round sau).
\end{itemize}
\end{itemize}

\item \textbf{Kết thúc game}
Giữ nguyên luật thắng thua của phiên bản gốc:
\begin{itemize}
\item Phe Thiện (Chó) thắng nếu có đường thông hợp lệ từ thẻ Start đến thẻ Goal chứa vàng.
\item Phe Ác (Mèo) thắng nếu hết bài trên tay tất cả người chơi mà vẫn chưa nối được đường đến vàng.
\end{itemize}
\item \textbf{Thẻ hành động mới (bổ sung thêm vào bộ bài)}
\begin{table}[H]
\centering
\begin{tabular}{|l|p{7cm}|p{5cm}|}
\hline
\textbf{Tên thẻ} & \textbf{Công dụng} & \textbf{Ghi chú} \\ \hline
Shield & Người chơi sẽ có khiên vĩnh viễn có thể đỡ được 1 đòn tấn công phá công cụ (Cart, Shovel, Hat) bởi thẻ Sabotage. & Có thể áp dụng cho bất cứ ai, không tích lũy \\ \hline
Binoculars & Khi sử dụng thẻ này, người chơi được nhìn thấy toàn bộ bản đồ từ lúc chơi thẻ cho đến hết Đêm hiện tại. & Có thể áp dụng cho bất cứ ai\\ \hline
Swap Goals & Trong Đêm, tráo vị trí ngẫu nhiên hoặc chỉ định của 3 thẻ Goal (khi chúng vẫn đang úp). & Thay đổi vị trí kho vàng thực sự \\ \hline
\end{tabular}
\caption{Các thẻ hành động mới}
\end{table}
\end{enumerate}